% ============================================================
%  Chapitre 2 — Données
% ============================================================

\chapter{Données}
\label{chap:donnees}

\section{Régions d'étude}

Deux régions agricoles des États-Unis ont été retenues, conformément au
protocole de Wang et al. \cite{wang2024} :

\begin{itemize}
  \item \textbf{Arkansas} — région à dominante rizicole et cotonnière dans le
        bassin du Mississippi. Cinq classes de cultures : Corn, Cotton, Rice,
        Soybeans, Others.
  \item \textbf{Californie} — région à horticulture intensive dans la Central
        Valley. Six classes : Rice, Alfalfa, Grapes, Almonds, Pistachios,
        Others.
\end{itemize}

Ces deux régions présentent des caractéristiques phénologiques contrastées,
ce qui permet d'évaluer la généralisation du modèle.

\section{Acquisition via Google Earth Engine}

Les données Sentinel-2 ont été acquises via la plateforme
\textbf{Google Earth Engine} (GEE), qui donne accès à l'archive complète des
images satellitaires sans téléchargement local.

Le pipeline d'acquisition comprend les étapes suivantes :
\begin{enumerate}
  \item Sélection de la collection \texttt{COPERNICUS/S2\_SR\_HARMONIZED}
        (Surface Reflectance, harmonisée).
  \item Filtrage par région d'intérêt (polygone délimitant l'État) et par
        période (saison agricole 2022).
  \item Masquage des nuages via la bande QA60 et le masque \texttt{SCL}.
  \item Rééchantillonnage temporel en 36 composites mensuels ou bi-mensuels
        par interpolation linéaire des valeurs manquantes.
  \item Sélection des 10 bandes spectrales retenues : B2, B3, B4, B5, B6,
        B7, B8, B8A, B11, B12.
  \item Export en format GeoTIFF, puis conversion en tableaux NumPy.
\end{enumerate}

Les vérités terrain (labels) proviennent du programme
\textbf{Cropland Data Layer} (CDL) de l'USDA, reclassifiées selon les
catégories de l'article.

\section{Zones d'échantillonnage}

Les pixels ont été échantillonnés dans des \textbf{zones dispersées} sur
l'ensemble de la région, afin d'éviter l'autocorrélation spatiale entre les
ensembles d'entraînement et de test — un biais courant lorsque les zones
d'entraînement et de test sont adjacentes. Cette stratégie s'écarte des
zones concentrées initialement envisagées et améliore la robustesse de
l'évaluation.

La répartition des échantillons est la suivante :

\begin{table}[H]
\centering
\caption{Répartition des échantillons par région et par split}
\label{tab:splits}
\begin{tabular}{lrrr}
\toprule
\textbf{Région} & \textbf{Train} & \textbf{Val} & \textbf{Test} \\
\midrule
Arkansas    & 1\,200 & 300 & 8\,523 \\
California  & 1\,440 & 360 & 8\,226 \\
\bottomrule
\end{tabular}
\end{table}

Le test set est volontairement large pour obtenir des métriques stables.

\section{Format des données}

Chaque pixel est représenté par :
\begin{itemize}
  \item \textbf{Input 1} (\texttt{input1.npy}) : tenseur de forme $(N, 10, 36)$,
        type \texttt{float32} — valeurs de réflectance normalisées.
  \item \textbf{Input 2} (\texttt{input2.npy}) : tenseur de forme $(N, 36)$,
        type \texttt{float32} — masque binaire des observations manquantes
        (1 = observation valide, 0 = nuage ou donnée absente).
  \item \textbf{Labels} (\texttt{labels.npy}) : vecteur de forme $(N,)$,
        type \texttt{int64} — indice de classe (0-indexé).
\end{itemize}

\section{Exploration statistique}

\subsection{Distribution des classes}

La distribution des classes dans le jeu d'entraînement reflète la
prédominance de certaines cultures dans chaque région. En Arkansas, la
classe Soybeans représente la majorité des pixels, tandis qu'en Californie,
Alfalfa et Grapes dominent.

% \begin{figure}[H]
%   \centering
%   \includegraphics[width=0.8\textwidth]{figures/class_distribution.png}
%   \caption{Distribution des classes dans le jeu d'entraînement.}
%   \label{fig:class_dist}
% \end{figure}

\subsection{Séries temporelles spectrales}

Les séries temporelles moyennes par classe montrent des profils phénologiques
distincts. Les cultures à cycle court (Rice en Arkansas) présentent un pic de
NDVI net entre les mois de juin et août, tandis que les cultures pérennes
(Alfalfa en Californie) affichent un signal élevé sur toute la saison.

\subsection{Données manquantes}

Le masque \texttt{input2} indique les pas de temps affectés par la couverture
nuageuse. En Arkansas, la couverture nuageuse estivale génère des lacunes sur
les mois de juillet-août. Le module ALPE (\textit{Adaptive Learnable
Positional Encoding}) de MCTNet a été spécifiquement conçu pour traiter ces
observations manquantes (cf. Section~\ref{sec:alpe}).

\section{Preprocessing}

\subsection{Normalisation}

Les valeurs de réflectance sont normalisées par bande à l'échelle du jeu
d'entraînement (mean/std). Cette normalisation est calculée uniquement sur
le train set et appliquée à val et test pour éviter la fuite de données.

\subsection{Gestion des valeurs manquantes}

Les positions masquées (masque = 0) conservent leur valeur d'interpolation
dans \texttt{input1}, mais le masque est explicitement fourni en entrée au
modèle (\texttt{input2}). Le Transformer intègre ce masque via ALPE pour
moduler la pondération de l'encodage positionnel selon la fiabilité de
l'observation.

\subsection{Découpage train/val/test}

Le découpage est effectué au niveau spatial avant l'échantillonnage des
pixels, de sorte qu'aucune zone géographique ne soit présente dans plusieurs
splits. La graine aléatoire est fixée à 42 pour la reproductibilité.

\cleardoublepage
